% !TeX root = 01_intro_ml_polytech.tex
% !TeX encoding = UTF-8
% !TeX spellcheck = en_GB

\begin{frame}{Some Introductions}{Let us talk about the course\dots}

    \begin{center}
        \huge
        Welcome to \\ \textbf{\enquote{Introduction to Machine Learning}}!
    \end{center}

    \vfill

    \begin{tabular}{@{}ll@{}}
        \textbf{Instructor:} & Riccardo Finotello -- \href{mailto:riccardo.finotello@cea.fr}{riccardo.finotello@cea.fr}                                                                                                                                                                                                                                             \\
                             & \emph{Ingénieur chercheur} {\tiny (uncertainty quantification, theory methods, material science)}                                                                                                                                                                                                                                    \\
                             & \emph{Laboratoire Intelligence Artificielle et Science des Données (LIAD)}                                                                                                                                                                                                                                                           \\[1em]
                             & \href{https://arxiv.org/a/finotello_r_1.html}{\includegraphics[height=12pt]{01_slides/img/arxiv.pdf}} ~~~ \href{https://github.com/thesfinox}{\includegraphics[height=12pt]{01_slides/img/github.pdf}} ~~~ \href{https://uranie.cea.fr/presentation/team/riccardofinotello}{\includegraphics[height=12pt]{01_slides/img/uranie.png}} \\[1em]
        \textbf{Slides:}     & \href{https://github.com/thesfinox/polytechparis-intro-to-ml}{https://github.com/thesfinox/polytechparis-intro-to-ml}
    \end{tabular}

    \begin{tikzpicture}[overlay]
        \node[anchor=east, xshift=-5em, yshift=2em] at (current page.south east) {\includegraphics[width=5em]{01_slides/img/slides_qr.png}};
    \end{tikzpicture}
\end{frame}

\begin{frame}{Some Introductions}{Syllabus}
    \vfill
    \begin{tabular}{@{}lp{35em}@{}}
        \textbf{Module 1} & \textbf{Motivation} for ML in physics and materials science -- \textbf{Python} project setup -- \textbf{Statistics foundations} brush up -- \textbf{Hands on}: Python basics and first data exploration                                                                                                                                                                                      \\ \hline
        \textbf{Module 2} & \textbf{ML best practices} (train/validation/test splits, feature scaling) -- \textbf{Optimisation perspective} (loss functions, constrained optimisation, regularisation). -- \textbf{Linear regression} (coefficient interpretations, metrics) -- \textbf{Logistic regression} (decision boundaries, classification metrics) -- \textbf{Hands on}: data splits, scaling, and linear models \\ \hline
        \textbf{Module 3} & \textbf{Support Vector Machines} (maximum-margin classification, kernel trick) -- \textbf{Decision trees} (ensemble principle) -- \textbf{Principal Components Analysis} (dimensionality reduction) -- \textbf{Clustering} (K-means, DBSCAN) -- \textbf{Hands on}: advanced exploration and algorithms                                                                                       \\ \hline
        \textbf{Module 4} & \textbf{Feed-forward networks} (perceptron) -- \textbf{Gradient Descent and Backpropagation} -- \textbf{Hands on}: PyTorch in practice and advanced ML
    \end{tabular}
    \vfill
\end{frame}

\begin{frame}{Some Introductions}{Finals}
    \vfill
    \begin{tabular}{@{}lp{35em}@{}}
        \textbf{Modality} & 1 written examination (on paper) on theoretical aspects of ML + interpretation of code (coefficient 100\%)                                                                          \\ \hline
        \textbf{When}     & TBD (based on your availability -- possible end 2026 or beginning 2027, but quick decision)                                                                                         \\ \hline
        \textbf{What}     & Entire syllabus (roughly 18/20) + bonus original (out-of-the-box thinking) questions                                                                                                \\ \hline
        \textbf{How}      & Course is taught in French, slides in English to help with language (ask questions in French, English, or Italian, if you prefer)                                                   \\ \hline
        \textbf{Rules}    & - AI is a tool to enhance learning (\emph{forbidden during finals}): use it at your discretion during lessons, avoid it when studying autonomously, use it to refine your knowledge \\
                          & - 2 hours for the finals (no notes allowed)                                                                                                                                         \\
                          & - No question is ever stupid: best to ask, rather than failing an interview                                                                                                         \\
                          & - Customary pause (\textasciitilde 15 min) between sessions, if agreed on
    \end{tabular}
    \vfill

    \usebeamerfont{footline}
    \usebeamercolor[fg]{footline}
    \tiny Failing to abide to these rules will result in immediate disqualification and report to the coordinator.
    \strut
\end{frame}
